\section{Wake, circulation, and non-equilibrium hypotheses}

The original ITSM described baryonic matter as exciting a torsional or
acoustic wake in a rotational medium.  The static equation
Eq.~\eqref{eq:periodic-p3} has no independent memory: for specified boundary
data it tracks the instantaneous source.  A persistent, advected or detached
wake therefore requires additional dynamics.

A schematic relaxation equation is
\begin{equation}
 \tau_W U^\mu\nabla_\mu W+W
 =\mathcal F[\psi,\rho_b,r,Q_{\rm syn}],
 \label{eq:wake-schematic}
\end{equation}
or, if wave propagation is essential, a hyperbolic field equation with a
second time derivative.  Equation~\eqref{eq:wake-schematic} is not adopted as a
law.  It identifies the quantities that WAK-001 must derive or bound.
Its first-order form is motivated only as a causal-relaxation template of the
Israel--Stewart type~\cite{israelstewart1979}; it is not a derivation of an
ITSM constitutive law.

A viable wake sector must satisfy all of the following:
\begin{enumerate}
\item a well-posed initial-value problem;
\item causal characteristic speeds in the physical frame;
\item a positive or bounded energy accounting;
\item a controlled static limit reproducing the IR field equation;
\item explicit momentum exchange with matter and the reservoir;
\item decay or transport laws rather than indefinite unsupported memory.
\end{enumerate}

Similarly, circulation is represented naturally by the phase in
Eq.~\eqref{eq:complex-order},
\begin{equation}
 \oint \nabla_i\Theta\,\dd x^i=2\pi n,
 \qquad n\in\mathbb Z,
\end{equation}
where zeros of $\rho$ permit defect cores.  The relationship between these
microscopic winding sectors and a global toroidal current is the VOR-001 gate.
No black-hole, lensing or galactic coefficient follows from the winding number
without further dynamics.

\subsection{3D multi-fluid cluster collision and causal wave propagation}

To evaluate whether causal retardation in the scale-compensator field $\psi(t, \mathbf{x})$ can separate gravitational lensing centroids from collisional gas during supersonic mergers (the Bullet Cluster 1E~0657-56 problem), gate WAK-001 formulates and solves the damped causal wave equation:
\begin{equation}
 \left(\frac{1}{c_s^2}\frac{\partial^2}{\partial t^2} + \frac{1}{\tau_W}\frac{\partial}{\partial t} - \nabla^2\right) \psi(t, \mathbf{x}) = 4\pi G V \left[\rho_{\rm gas}(t, \mathbf{x}) + \rho_{\rm stars}(t, \mathbf{x})\right],
 \label{eq:causal-wave-3d}
\end{equation}
where $c_s = c/\sqrt{3} \approx 176.9\,\mathrm{kpc/Myr}$ and $\tau_W \approx 15\,\mathrm{Myr}$.

Solving Eq.~\eqref{eq:causal-wave-3d} across a 3D multi-fluid collision ($v_{\rm infall} = 4500\,\mathrm{km/s}$) via an exact Fourier-space matrix exponential propagator demonstrates that:
\begin{enumerate}
 \item \textbf{Transient Centroid Separation:} At the supersonic transit epoch ($t = 25\,\mathrm{Myr}$), the dynamic phase lag of the scalar response produces an effective lensing centroid offset of $\Delta x_{\rm lens-gas} \approx 20.91\,\mathrm{kpc}$, consistent with observed weak-lensing separations ($\Delta x_{\rm obs} \approx 20\text{--}30\,\mathrm{kpc}$).
 \item \textbf{Late-Time Relaxation:} In the absence of non-linear hydrodynamic screening or vortex phase disruption of the shocked gas, the long-term scalar wake relaxes back toward the dominant baryonic gas mass at $t = 35\,\mathrm{Myr}$ ($\Delta x \to -7.20\,\mathrm{kpc}$).
\end{enumerate}
Consequently, WAK-001 proves that linear causal retardation is kinematically capable of producing transient post-collision lensing offsets, but full non-linear Landau phase screening is required for permanent post-merger decoupling. It is classified as an \textit{Exploratory Kinematic Scaffold}.
